\section{Strict cones, finite class sectors, and cyclotomic rank}
\label{sec:cubic-frame-mu-criterion}

\paragraph{Recovery revision.}
This manuscript is a new recovery revision written on 22 November 2026.
The earlier AC-G74 text was absent from the restored checkpoint and the
identified private stores.  The proofs below reconstruct the retained
mathematical statements, with fresh literal-source checks.  No identity
with the missing original bytes is asserted.  The local-map correction
and the arbitrary-rank specialization example are included explicitly.

All cohomology is continuous.  The coefficient prime \(p\) is odd and
\(k\) is a finite field of characteristic \(p\).  A local map in a cone
means one specified homomorphism, counted with its multiplicity in the
index set.  It need not mean one distinct prime.

\begin{proposition}[AC-G74-cone: the strict kernel and its local correction]
\label{prop:ACG74cone}
Let \(G\) be profinite, let \(V\) be a finite-dimensional continuous
\(k\)-representation, and let
\(\operatorname{loc}_i:G_i\longrightarrow G\), \(i\in J\), be a finite
family of continuous homomorphisms from profinite groups.  Put
\[
 H^1_{\mathrm{str}}(G,V)=
 \ker\left(H^1(G,V)\longrightarrow
                  \bigoplus_{i\in J}H^1(G_i,V)\right).
\]
For the source's cone with zero local subspaces there is an exact sequence
\[
 0\longrightarrow
 \operatorname{coker}\left(V^G\longrightarrow
                   \bigoplus_{i\in J}V^{G_i}\right)
 \longrightarrow H^1_C(\operatorname{loc},V,0)
 \longrightarrow H^1_{\mathrm{str}}(G,V)
 \longrightarrow0.                                             \tag{G74.1}
\]
In particular its correction to the strict dimension is at most
\(|J|\dim_k V\); it vanishes if every \(V^{G_i}=0\).
For the full local subspaces,
\[
                    H^1_C(\operatorname{loc},V,\mathrm{top})
                         \simeq H^1(G,V).                       \tag{G74.2}
\]

Let \(L/F\) be a finite Galois extension of number fields of degree
prime to \(p\), with group \(\Delta\), trivializing \(V\).
Require a class to be unramified away from \(p\) after restriction to
\(G_L\), and zero at every full decomposition group above \(p\).
Then this strict subgroup is canonically
\[
 \operatorname{Hom}_{k[\Delta]}
 \left(\operatorname{Cl}(\mathcal O_L[1/p])_{(p)}/p
                          \otimes_{\mathbf F_p}k,V\right).
                                                               \tag{G74.3}
\]
A finite extension of degree prime to \(p\) may be used to test each
local zero condition.  If \(V=k(\theta)\), where
\(\theta:\Delta\longrightarrow k^\times\), the dimension in (G74.3)
is the dimension of the \(\theta\)-isotypic part, with idempotent
\[
             e_\theta=\frac1{|\Delta|}
                      \sum_{\delta\in\Delta}\theta(\delta)^{-1}\delta .
                                                               \tag{G74.4}
\]
There is no inversion of \(\theta\) in this dimension formula.
\end{proposition}

\begin{proof}
For zero local subspaces the literal source cocycles are pairs
\((z,(t_i))\) satisfying
\[
 z(gh)=z(g)+g z(h),\qquad
 z(\operatorname{loc}_i(s))=\operatorname{loc}_i(s)t_i-t_i.
\]
The boundaries are the pairs
\((g\mapsto gv-v,(v)_i)\).
Projection to \([z]\) maps onto the strict kernel: its local
vanishing is precisely the existence of the displayed primitives
\(t_i\).  If \([z]=0\), subtract a global boundary and obtain a tuple
in \(\bigoplus_i V^{G_i}\).  Two such tuples differ by a boundary
exactly when they differ by a diagonal element of \(V^G\).
This proves (G74.1), including its dimension bound.
For full local subspaces there is no constraint on \((t_i)\),
and every tuple with \(z=0\) is a boundary; projection gives
(G74.2).

Since \(|\Delta|\) is invertible in \(k\), inflation--restriction
identifies \(H^1(G_F,V)\) with
\(H^1(G_L,V)^\Delta\).  The latter is
\(\operatorname{Hom}_{\mathrm{cont}}(G_L,V)^\Delta\).
The required homomorphisms kill all away-\(p\) inertia and all
\(p\)-decomposition groups.  Global class field theory identifies
the resulting maximal abelian pro-\(p\) quotient with
\(\operatorname{Cl}(\mathcal O_L[1/p])_{(p)}\).
Real-place issues disappear for odd \(p\).
This proves (G74.3).  Restriction to an extension of degree prime
to \(p\) is injective on local \(H^1\), by
restriction--corestriction, so it gives the same zero test.

The equivariance convention is
\[
                 f(\delta a\delta^{-1})=\theta(\delta)f(a).
                                                               \tag{G74.5}
\]
Thus \(f\) is supported on \(e_\theta\), rather than
\(e_{\theta^{-1}}\).  The operator (G74.4) projects onto the
subspace on which each \(\delta\) acts by \(\theta(\delta)\);
on this subspace every \(k\)-linear functional is equivariant
to \(k(\theta)\).  This proves the final assertion.
\end{proof}

\begin{theorem}[AC-G74: exact finite-layer control and a finite rank certificate]
\label{thm:ACG74}
Let \(F_\infty/F\) be the cyclotomic \(\mathbf Z_p\)-extension,
and let \(L/F\) be finite Galois of degree prime to \(p\), with
group \(\Delta\).  Put \(L_\infty=LF_\infty\).
After replacing \(F,L\) by corresponding finite layers, assume
that every prime above \(p\) is totally ramified in \(L_\infty/L\).
Write
\[
 \Gamma=\operatorname{Gal}(L_\infty/L)\simeq\mathbf Z_p,\qquad
 \Gamma_n=\Gamma^{p^n},\qquad L_n=L_\infty^{\Gamma_n},\qquad
 A_n=\operatorname{Cl}(\mathcal O_{L_n}[1/p])_{(p)} .
\]
Let \(\Sigma_p\) be the finite \(\Delta\)-set of primes of \(L\)
above \(p\), identified with the primes at every higher layer.
Let \(\mathcal M/L_\infty\) be the maximal abelian pro-\(p\)
extension unramified at every finite prime and split completely
at every prime above \(p\), and set
\(X'=\operatorname{Gal}(\mathcal M/L_\infty)\).

There is a \(\Delta\)-equivariant exact sequence
\[
 \mathbf Z_p[\Sigma_p]^0
       \longrightarrow (X')_{\Gamma_n}
       \longrightarrow A_n\longrightarrow0,\qquad
 \mathbf Z_p[\Sigma_p]^0=
       \ker\bigl(\mathbf Z_p[\Sigma_p]\xrightarrow{\mathrm{aug}}
                                                   \mathbf Z_p\bigr).
                                                               \tag{G74.6}
\]
No injectivity of its left map is asserted.  In particular the
direction of the right map is from coinvariants onto the
finite-layer class group.

Choose a generator \(\gamma\) of \(\Gamma\), put \(T=\gamma-1\),
and let \(\theta:\Delta\to k^\times\).  Define
\[
 \begin{split}
 M_\theta&=e_\theta\bigl((X'/p)\otimes_{\mathbf F_p}k\bigr),\\
 a_{n,\theta}&=\dim_k e_\theta
                      \bigl((A_n/p)\otimes_{\mathbf F_p}k\bigr),\\
 b_\theta&=\dim_k e_\theta\bigl(k[\Sigma_p]^0\bigr).
 \end{split}                                                    \tag{G74.7}
\]
Then \(M_\theta\) is a finitely generated \(k[[T]]\)-module.
There are integers \(r_\theta,c_\theta\geq0\), independent of \(n\),
such that
\[
 \begin{split}
 r_\theta p^n
 &\leq \dim_k M_\theta/T^{p^n}M_\theta
 \leq a_{n,\theta}+b_\theta,\\
 a_{n,\theta}&\leq r_\theta p^n+c_\theta.
 \end{split}                                                    \tag{G74.8}
\]
Consequently every single finite level \(m\) supplies the certificate
\[
 r_\theta\leq
 \left\lfloor\frac{a_{m,\theta}+b_\theta}{p^m}\right\rfloor .
                                                               \tag{G74.9}
\]
If \(X'\) is finitely generated over \(\mathbf Z_p\), then
\(r_\theta=0\).  No such hypothesis on the actual cubic
coefficient splitting field is supplied by this theorem.
\end{theorem}

\begin{proof}
\proofstep{The two commuting actions and the decomposition sections.}
Since \([L:F]\) is prime to \(p\),
\(L\cap F_\infty=F\), and
\(\operatorname{Gal}(L_\infty/F)=\Delta\times\Gamma\).
The defining maximality of \(\mathcal M\) makes it Galois over \(F\).
The \(\Delta\)- and \(\Gamma\)-actions on its abelian subgroup \(X'\)
are well defined, independent of lifts, and commute.

Fix \(n\), and put \(H_n=\operatorname{Gal}(\mathcal M/L_n)\).
For a prime of \(\mathcal M\) above \(v\in\Sigma_p\), its
decomposition subgroup maps isomorphically onto \(\Gamma_n\):
the kernel is trivial because \(\mathcal M/L_\infty\) is
\(p\)-split, and the image is all of \(\Gamma_n\) because
\(L_\infty/L_n\) is totally ramified at that prime.
Any one such subgroup splits
\[
                 1\longrightarrow X'\longrightarrow H_n
                         \longrightarrow\Gamma_n\longrightarrow1.
\]
Abelianizing the resulting semidirect product gives
\[
       0\longrightarrow (X')_{\Gamma_n}
         \longrightarrow H_n^{\mathrm{ab}}
         \longrightarrow\Gamma_n\longrightarrow0.
                                                               \tag{G74.10}
\]
Commutators and coinvariant images here are closed.  Equivalently,
\((\gamma^{p^n}-1)X'\) is compact and hence closed.

For every \(v\), the decomposition section in
\(H_n^{\mathrm{ab}}\) is independent of the chosen prime above it,
since changing the prime conjugates its subgroup inside \(H_n\).
Denote it by \(s_v\).
Writing \(\gamma_n=\gamma^{p^n}\), define
\[
 \beta_n:\mathbf Z_p[\Sigma_p]\longrightarrow H_n^{\mathrm{ab}},
                 \qquad [v]\longmapsto s_v(\gamma_n).
                                                               \tag{G74.11}
\]
Its composite with \(\Gamma_n\simeq\mathbf Z_p\) is the augmentation.
The \(\Delta\)-action on \(H_n^{\mathrm{ab}}\) is canonical even
when a lift of \(\Delta\) to the larger Galois group is chosen
only up to an inner automorphism of \(H_n\).
As \(\Delta\) commutes with \(\Gamma_n\) and permutes the primes,
\(\beta_n\) is \(\Delta\)-equivariant.

\proofstep{Both class-field directions.}
The quotient of \(H_n^{\mathrm{ab}}\) by
\(\beta_n(\mathbf Z_p[\Sigma_p])\) kills the full decomposition
groups at \(p\).  It is unramified elsewhere: both
\(L_\infty/L_n\) and \(\mathcal M/L_\infty\) are unramified
away from \(p\).  It is therefore a quotient of \(A_n\).

For the converse, let \(\mathcal A_n/L_n\) be the maximal
unramified abelian \(p\)-extension in which all \(p\)-primes split;
its Galois group is \(A_n\), by class field theory.
The intersection \(\mathcal A_n\cap L_\infty\) is \(L_n\).
Indeed a \(p\)-prime splits completely in the first extension
and is totally ramified in the second, so the intersection
has trivial decomposition group and full decomposition group
simultaneously.
The compositum \(\mathcal A_n L_\infty/L_\infty\) is abelian
pro-\(p\), unramified, and \(p\)-split, hence is contained in
\(\mathcal M\).  Thus \(A_n\) is the full quotient just described,
not merely an unspecified class-field factor.

Every element of \(H_n^{\mathrm{ab}}\) can be made to lie in
\((X')_{\Gamma_n}\) by subtracting the image of one section.
The intersection
\[
 (X')_{\Gamma_n}\cap\beta_n(\mathbf Z_p[\Sigma_p])
                 =\beta_n(\mathbf Z_p[\Sigma_p]^0)
\]
follows by applying the projection in (G74.10).
This proves (G74.6), with its stated equivariance.

\proofstep{Finite generation, without a hidden Iwasawa-module premise.}
Reduce (G74.6) modulo \(p\), extend scalars to \(k\), and apply
the exact idempotent \(e_\theta\).
At \(n=0\) this shows that \(M_\theta/TM_\theta\) is finite-dimensional,
because \(A_0\) is finite and the augmentation module has finite rank.

For completeness, compact Nakayama applies here as follows.
Choose finitely many lifts of a basis of this quotient.
Their \(k[[T]]\)-span \(M_0\) is the compact image of a finite free
\(k[[T]]\)-module, hence is closed.
The compact quotient \(Q=M_\theta/M_0\) satisfies \(Q=TQ\).
There is a basis of neighborhoods of zero consisting of open
\(\Gamma\)-stable subspaces: continuity and compactness of the
\(\Gamma\)-action give a common open neighborhood inside all
translates of a given open subspace, and one can then take its
stable span.
On every resulting finite quotient the pro-\(p\) group \(\Gamma\)
has finite \(p\)-power image, so \(T\) is nilpotent.
A finite quotient of \(Q=TQ\) is therefore zero.
Since these finite quotients separate points of \(Q\), \(Q=0\).
This proves the claimed finite generation.

The structure theorem over the discrete valuation ring \(k[[T]]\)
now gives
\[
 M_\theta\simeq k[[T]]^{\,r_\theta}
              \oplus\bigoplus_{i=1}^{h}k[[T]]/(T^{a_i}),
              \qquad c_\theta=\sum_{i=1}^{h}a_i ,
                                                               \tag{G74.12}
\]
where the list may be empty.
In characteristic \(p\), \(\gamma^{p^n}-1=T^{p^n}\), so
\[
 \dim_k M_\theta/T^{p^n}M_\theta
       =r_\theta p^n+\sum_i\min(a_i,p^n).
                                                               \tag{G74.13}
\]
The reduction of (G74.6) gives a surjection from this quotient
onto the \(\theta\)-part of \(A_n/p\), with kernel dimension at
most \(b_\theta\).  Equations (G74.8) and (G74.9) follow.
If \(X'\) is \(\mathbf Z_p\)-finite, \(X'/p\) is finite, and the
free \(k[[T]]\)-rank is zero.
\end{proof}

\paragraph{What the estimates do and do not supply for a frame.}
For a fixed tower and a fixed finite collection of strict characters,
(G74.8) bounds every strict dimension by its \(k[[T]]\)-rank times
\(p^n\), plus a constant independent of \(n\).
The same is true of a fixed sum of sectors.
Equation (G74.1) adds a uniform constant only when the number of
specified local maps is uniformly bounded.  This holds, for example,
for one map per \(p\)-prime after the above reindexing, or for a fixed
number of repeats of each map.  Arbitrarily many duplicate maps are
not covered by that assertion.  If all local invariant spaces vanish,
there is no correction at all.

In the literal source, \path{Frame.Ok.f3} makes the local values
of \(\Psi^2\) nontrivial.  Since the local maps land in
\(\ker\delta\), it kills the local invariants for
\(\Psi^{-1}\), \(\Psi^{-1}\delta\), and
\(\Psi^{-2}\delta\).  The two nontriviality clauses are
\[
 H^1_C(\operatorname{loc},k(\Psi^{-1}\delta),0)\ne0,\qquad
 H^1_C(\operatorname{loc},k(\Psi^{-2}\delta),0)\ne0.
                                                               \tag{G74.14}
\]
They are independent requirements; an upper-rank estimate does
not prove either one.  The source bounds
\path{Frame.Ok.b1p} and \path{Frame.Ok.b1} concern, respectively,
the second cone in (G74.14) and the
\(\Psi^{-1}\)-cone on the CM kernel with its doubled local maps.
The full-local clause \path{Frame.Ok.h1} concerns trivial
coefficients and \(\mathrm{top}\), so it is (G74.2), not an
additional strict character sector.  Its exact arithmetic
comparison and the dihedral-host bound are given separately
in AC-G76.

For example, if the required strict bounds have slopes at most
\(R\), an independently verified full-local bound fits the same
choice \(C_{0,n}=Rp^n+C\), and \(d_n=d_0p^n\), then
\(d_0>3R\) gives \(3C_{0,n}+7<d_n\) for sufficiently large \(n\).
If every source local degree also grows by \(p^n\), its condition
\(d_{v,n}>5\) eventually holds.  Formula (G74.9) gives a finite
way to certify the needed slopes, but the required finite
class-sector values and the strict inequality are not computed
here.  Nor is an actual replacement of the single-power source
frame by two different cubic adjacent coefficients constructed.

\begin{proposition}[AC-G74-specialization: zero analytic \(\mu\) does not bound the specialized rank]
\label{prop:ACG74specialization}
Let \(\mathcal O\) be a discrete valuation ring with uniformizer
\(\pi\) and residue field \(k\).  Put
\(\Lambda_2=\mathcal O[[T_+,T_-]]\).  For every integer \(N\geq1\)
there is a finitely generated torsion \(\Lambda_2\)-module whose
full characteristic ideal and anticyclotomic specialization both
have \(\mu=0\), while its cyclotomic specialization modulo \(\pi\)
has \(k[[T_+]]\)-rank \(N\).
Thus these two analytic \(\mu=0\) assertions alone yield no
numerical bound on that specialized rank.
\end{proposition}

\begin{proof}
Take
\[
          \mathcal N_N=
               \bigl(\Lambda_2/(T_-+\pi)\bigr)^{\oplus N}.
                                                               \tag{G74.15}
\]
Its characteristic ideal is \(((T_-+\pi)^N)\), which is not
divisible by \(\pi\).  Specializing \(T_+=0\) gives
\[
       \bigl(\mathcal O[[T_-]]/(T_-+\pi)\bigr)^{\oplus N},
\]
whose characteristic ideal likewise has no \(\pi\)-factor.
On the other hand,
\[
        \mathcal N_N/T_-\mathcal N_N
                          \simeq k[[T_+]]^{\,N}.
                                                               \tag{G74.16}
\]
This module is already killed by \(\pi\), and its indicated rank
is \(N\).  Since \(N\) is arbitrary, the asserted formal
implication fails.  This is a module counterexample, not a
claim that a particular arithmetic Selmer module has this form.
\end{proof}

\paragraph{Exact source and tower scope.}
The cone definitions and every frame clause just quoted are read
literally in \path{ContinuousCoboundaries2.lean}, at pinned commit
\path{7cbf7a0c32fca1badf05931879a37c3f8a2b03b4}, under
\path{Development/AutomorphicArithmetic/AutomorphicVarieties/FrobeniusSmoothVectors/}.
The applied definitions are \path{Z1C},
\path{continuousCoboundaries1}, \path{H1C}, \path{Frame.Ok},
and \path{Frame.Big}; their complete containing file was read
for this recovery revision.
The source's \path{HostFrameEx.lean} imposes
\path{GalLink.ab}, so a nonabelian dihedral host is not a literal
instance of that abelian-field hypothesis without an additional
proved replacement.

The class-field statements used are Milne,
\emph{Class Field Theory}, version 4.03, V.3.5--3.7,
Example V.3.9, and the \(S\)-class formulation in Exercise V.3.15.
These identify the unramified class field and its quotient split
at the specified primes.  They are imported class-field theorems;
no independent proof of the entire reciprocity or existence
theorem is claimed.  Sharifi, \emph{Iwasawa Theory},
Propositions 3.1.9 and 3.1.11, supplies that the cyclotomic tower
is unramified away from \(p\) and that the \(p\)-primes are
eventually totally ramified.  The complete printed proofs of
these two applied propositions were read.  The finite-layer
comparison and compact-module arguments above are proved here.

Hsieh's \emph{Iwasawa main conjectures for CM fields},
\(\S8.2.1\), printed pages 92--93, defines its
\(\Sigma\)-ordinary condition using inertia at the
\(\Sigma^c\)-places; it does not impose zero restriction to every
\(p\)-adic decomposition group.  Its cyclotomic and
anticyclotomic variables are set up in \(\S8.1\).
The literal local condition and the specialization map therefore
require a proved comparison before that theorem can supply the
strict sectors used here.  (G74.15)--(G74.16) explains a separate
algebraic obstruction to inferring the needed numerical
cyclotomic rank merely from characteristic-ideal \(\mu=0\).
It refutes no published arithmetic theorem and supplies no
actual class-sector inequality.

\paragraph{Geometric dictionary and remaining scope.}
The class groups describe abelian finite covers unramified at
finite primes; quotienting by the classes of the \(p\)-primes
forces those primes to split.  Restricting coefficients and
taking \(\Delta\)-isotypic parts is the corresponding operation
on finite local systems.  This is within the already audited
SGA~1, Expos\'e~V, \(\S\S4,7\), dictionary; no new archival
source read or new SGA theorem is claimed.
The result is a finite-certificate criterion for strict
cohomology.  Actual class-sector computations, positivity,
all remaining frame clauses, the two-adjacent-character
replacement, and prescribed Hecke-curve support remain distinct
requirements.

